\documentclass[12pt,a4paper]{article}
\usepackage{fullpage}
\usepackage{setspace}
\setstretch{1.2}
\usepackage[backend=biber,style=authoryear]{biblatex}
\addbibresource{../ref.bib}
\usepackage{hyperref}
\hypersetup{colorlinks=true, allcolors=blue}

\newcommand{\pkg}[1]{#1}
\newcommand{\proglang}[1]{\textsf{#1}}

\begin{document}

\noindent \textbf{Responses to the editorial comments on the previous submission
``RprobitB: Bayes Estimation of Choice Behavior Heterogeneity in R'' (JSS 4817),
which accompany the new submission ``RprobitB: Bayesian Probit Choice Models
with Preference Heterogeneity in R''}\\

We thank the editors for their comments on our previous submission. We have
since rewritten the manuscript and revised the \proglang{R} package
\texttt{\{RprobitB\}} \parencite{Oelschlaeger2026d}. Below, we quote each
comment and state how it is addressed. The software changes are listed at the
end of this document.
\\

\noindent \textbf{Responses to the editorial report}:

\begin{enumerate}
    \item Comment: The introduction indicates that a number of other software
    packages provide related functionality. It would seem that the specific
    added functionality of RprobitB compared to other packages is ``LC'' and
    ``LC update'', i.e., weight-based class updates as well as a Dirichlet
    process mixture approach. However, this aspect of the implementation is not
    sufficiently well discussed or presented using applications to be
    convincing.
    \item[] Response: We now state the contribution of the package more clearly
    and have improved the comparison with competing packages. We also discuss
    the latent class models in more detail and have reworked their application.

    Section~1 describes how competing packages model preference heterogeneity.
    The selection is guided by two central aspects of \texttt{\{RprobitB\}},
    namely probit models and preference heterogeneity. For preference
    heterogeneity, we include the established open-source packages
    \texttt{\{mlogit\}}, \texttt{\{gmnl\}}, \texttt{\{logitr\}},
    \texttt{\{apollo\}}, and \texttt{Biogeme}, which model it mostly via logit
    models. For probit models, we include \texttt{\{bayesm\}} and
    \texttt{\{MNP\}} for Bayesian estimation, \texttt{\{Rchoice\}} for random
    coefficients in binary and ordered models, and \texttt{\{Rprobit\}}, the
    frequentist counterpart of \texttt{\{RprobitB\}}. Table~1 shows for each of
    them whether it offers probit models, Bayesian estimation, random
    coefficients, latent classes, and inference on the number of classes.

    As Table~1 shows, \texttt{\{RprobitB\}} is, to the best of our knowledge,
    the only package that estimates multinomial probit models with random
    coefficients and correlated errors in a Bayesian framework, and the only one
    that infers the number of latent classes of such a model from the data;
    \texttt{\{bayesm\}} does so for logit models. The added functionality is not
    limited to the latent classes: for these probit models, the package offers
    normal and log-normal mixing distributions (Sections~2.4, 5.1, and 5.3),
    inference on individual deciders and conditional predictions (Sections~5.2
    and 7.1), and model comparison (Section~8).

    The latent class functionality is now discussed in detail in Sections~2.5,
    3.2, and 3.3 (see our responses to the next three comments). The application
    in Section~5.4 now uses the \texttt{Electricity} data, for which a single
    normal mixing distribution cannot represent a skewed or multimodal valuation
    of a local supplier. The Dirichlet process mixture and the sparse finite
    mixture infer the number of occupied classes, and
    \texttt{latent\_class\_diagnostics()} reports its posterior distribution.
    The weight-based heuristic is demonstrated on simulated data with a trimodal
    mixing distribution (Figure~5). Further applications cover a stated choice
    experiment on train trips (Sections~3.5, 4, 5.3, 7.3, and 8), the
    \texttt{Health} data for ordered responses (Sections~6.1 and 7.2), the
    \texttt{Game} data for ranked responses (Section~6.2), and Berserk decisions
    in an online chess tournament for conditional predictions (Section~7.1).
    \item Comment: The weighted-based class updates approach seem to have only
    been proposed by the authors in a conference publication and still lacks a
    sufficient statistical basis. It would seem that this is a rather ad-hoc
    approach for estimating the number of classes requiring various tuning
    parameters where no insights are provided for the reader how to choose or
    select them and what their impact is from a theoretical perspective
    regarding a constraint estimation or from a Bayesian point of view how this
    could be translated to a prior.
    \item[] Response: We agree that the weight-based heuristic lacks a formal
    statistical basis. It is therefore no longer the main contribution of the
    package, and the manuscript presents it as one of four treatments of the
    number of classes, which the argument \texttt{class\_update} of
    \texttt{fit()} selects (Sections~2.5 and 3.3): (i) a fixed number of classes
    (\texttt{"fixed"}); (ii) an overfitted finite mixture with a sparse
    Dirichlet prior on the class weights
    \parencite{Rousseau2011,MalsinerWalli2016,FruehwirthSchnatter2019}
    (\texttt{"sparse"}); (iii) a Dirichlet process mixture
    \parencite{Escobar1995}, sampled by Algorithm~8 of \textcite{Neal2000}
    (\texttt{"dirichlet\_process"}); and (iv) the heuristic of
    \textcite{Oelschlaeger2021} (\texttt{"weight\_based"}).

    From a Bayesian point of view, options (ii) and (iii) address the same task
    through a prior. In the sparse finite mixture, a sparse prior on the class
    weights empties the classes that the data do not require, similar to the
    removal of small classes by the heuristic, and the Dirichlet process mixture
    makes the number of occupied classes part of the posterior. We therefore
    recommend these options for inference on latent classes, and the heuristic
    primarily for the approximation of non-normal mixing distributions by a
    mixture of normal distributions \parencite{Oelschlaeger2026c}. We do not
    attach a prior interpretation to the heuristic. It changes the number of
    classes only during the warmup (Section~5.4). Afterwards, each chain keeps
    the number of classes it has reached. The retained draws thus come from an
    ordinary finite mixture with this number of classes and do not reflect the
    uncertainty about the number of classes, which options (ii) and (iii)
    capture.

    Section~5.4 states the main tuning parameters and their defaults, which
    worked well in our simulations: a class is removed if its weight falls below
    \texttt{epsmin = 0.01} and split if its weight exceeds
    \texttt{epsmax = 0.7}, two classes are merged if their means are closer than
    \texttt{deltamin = 0.1}, and at most one change occurs per
    \texttt{buffer = 50} iterations. Suitable values also depend on the specific
    application, which makes general recommendations difficult. Larger values of
    \texttt{epsmin} and \texttt{deltamin} tend to reduce the number of classes.
    Since only splits add classes, no class is added while all weights stay
    below \texttt{epsmax}, and a smaller \texttt{epsmax} tends to increase the
    number of classes. In a simulation with three equally weighted classes, the
    heuristic converges to three classes when starting from one and from ten
    initial classes (Figure~5).
    \item Comment: Related work on finite and infinite mixture models would
    indicate that in particular the Dirichlet parameter delta for finite
    mixtures and the precision parameter for DPMs would seem influential (see
    for example Fr\"uhwirth-Schnatter and Malsiner-Walli, 2019). In the current
    exposition it remains unclear how the model could be estimated if the number
    of classes should be a-priori specified to be filled and for example
    model-selection based on information criteria performed in the finite case,
    or a sparse finite mixture approach pursued or how the precision parameter
    is to be selected in DPMs. It would seem important to provide more
    theoretical as well as practical guidance on how to specify and apply the
    mixture approach in practice.
    \item[] Response: The manuscript now addresses the three approaches
    mentioned in the comment, both methodologically (Sections~2.5, 3.2, and 3.3)
    and in an application to the same data (Sections~5.4 and 8.2). As
    Section~3.2 and Table~4 state, the class weights of a finite mixture have a
    symmetric Dirichlet prior with concentration parameter $\delta$, which by
    default equals 1 for \texttt{"fixed"} and \texttt{"weight\_based"}, making
    the prior uniform on the simplex.

    With \texttt{class\_update = "fixed"}, the number of classes is specified a
    priori, and the sampler targets the posterior restricted to allocations
    without empty classes (Section~3.3). Model selection then compares fits with
    different numbers of classes: \texttt{AIC()} and \texttt{BIC()} use
    \texttt{logLik()}; \texttt{WAIC()} and \texttt{loo()} compute WAIC and
    PSIS-LOO at the decider level; and \texttt{bayes\_factor()} estimates Bayes
    factors (Section~8). For the \texttt{Electricity} data, PSIS-LOO favors two
    classes over one.

    With \texttt{class\_update = "sparse"}, the specified number of classes is
    an upper bound, and $\delta$ has a gamma hyperprior with shape~1 and
    rate~200, as in \textcite{FruehwirthSchnatter2019}. With six potential
    classes, the sparse finite mixture places a posterior probability of
    91~percent on two occupied classes (Section~5.4).

    With \texttt{class\_update = "dirichlet\_process"}, $\delta$ is the
    concentration, or precision, parameter of the process. The sampler updates
    it by the augmentation of \textcite{Escobar1995} under their gamma
    hyperprior with shape~2 and rate~4. For the 361 households, the prior
    expected number of occupied classes is about 3.9 at the prior mean
    $\delta = 0.5$ (Section~5.4), and the posterior places 47~percent on two
    classes and 53~percent on three or more. For the sparse finite mixture with
    six potential classes, the prior expected number of occupied classes is only
    about 1.2 at the prior mean $\delta = 0.005$. Under the matching of
    \textcite{FruehwirthSchnatter2019}, this corresponds to a Dirichlet process
    whose concentration has prior mean 0.03 instead of 0.5. With this matched
    hyperprior, the Dirichlet process places 87~percent on two classes, close to
    the 91~percent of the sparse finite mixture. The following code reproduces
    it from \texttt{electricity\_classes} of Section~5.4:
\begin{verbatim}
set.seed(1)
matched <- update(electricity_classes,
  class_update = "dirichlet_process",
  prior = list(class_concentration = c(shape = 1, rate = 200 / 6)))
subset(latent_class_diagnostics(matched)$occupancy, n_classes == 2)
#>   n_classes probability
#> 2         2   0.8686667
\end{verbatim}
    The difference between the two mixtures is thus mainly due to their
    hyperpriors, in line with the finding of \textcite{FruehwirthSchnatter2019}
    that the hyperprior influences the cluster solution far more than the choice
    between the two mixtures. In practice, we recommend checking the number of
    occupied classes under a second hyperprior, which can be set through
    \texttt{prior}.
    \item Comment: Label switching also seems to be only resolved using a
    constraint on the class weights. Given the large literature on resolving
    this issue for Bayesian mixture models and the failure of using only a
    constraint on one specific parameter in some cases this seems not
    satisfactory.
    \item[] Response: The constraint on the class weights has been removed.
    Following the literature on label switching
    \parencite[e.g.,][]{Stephens2000}, the sampler imposes no ordering
    constraint, and the retained draws are relabeled after sampling
    (Section~3.3). A reference allocation is selected by the least-squares
    criterion of \textcite{Dahl2006} with respect to the posterior co-clustering
    matrix, and each draw is permuted to agree with it by the equivalence
    classes representatives algorithm of \textcite{Papastamoulis2010}. The
    reference is then reset to the most frequent relabeled class of every
    decider until it is stable, and all class-specific parameters and the
    allocations are permuted accordingly. Since the relabeling is based on the
    allocations of all deciders, it does not require the classes to be separated
    in one particular parameter, whereas an ordering constraint does. The
    classes are then numbered by decreasing posterior mean weight. In the
    simulation of Section~3.4, whose three classes have equal weights and
    therefore cannot be identified by an ordering constraint on the weights, the
    relabeled class weights and means are close to their true values, with a
    rank-normalized $\hat{R}$ \parencite{Vehtari2021} of at most 1.01.
    \item Comment: We decided not to send the submission into full review
    because in its current form the chances for successful publication are
    deemed too low. The package provides a nice set of functionality making use
    of R's classes and methods system and building on infrastructure available,
    e.g., for formulas. However, it is not sufficiently well outlined how the
    package provides additional functionality compared to other packages and the
    potential additional functionality is not well enough theoretically
    explained, the statistical details and their implementation presented and
    their use illustrated in practice.
    \item[] Response: The package retains its S3 classes and methods and its
    formula interface. The additional functionality compared to other packages
    is outlined in Section~1 and Table~1 (see our response to the first
    comment). It is explained theoretically in Section~2, which defines the
    model class with random coefficients, latent classes, and ordered and ranked
    responses, together with the likelihood. The statistical details and their
    implementation are presented in Section~3: the function \texttt{fit()} and
    its arguments (Section~3.1 and Table~3), the prior distributions with their
    defaults (Section~3.2 and Table~4), the main blocks of the Gibbs sampler,
    the treatment of the number of classes, and the relabeling (Section~3.3). A
    simulation checks whether the sampler recovers the parameters of a known
    data generating process (Section~3.4). The use in practice is demonstrated
    in the applications listed in our response to the first comment, always with
    code illustrations, from the interpretation of the estimates (Section~4) to
    model evaluation (Section~8).
\end{enumerate}

\noindent \textbf{Responses to the specific comments}:

\begin{enumerate}
    \item Comment: The NAMESPACE contains
\begin{verbatim}
S3method(AIC,RprobitB_fit) S3method(BIC,RprobitB_fit)
\end{verbatim}
    raising the question why the logLik method does not return a logLik object
    where AIC and BIC methods are readily available for.
    \item[] Response: \texttt{logLik()} now returns a \texttt{logLik} object
    with the number of parameters and the number of observations as attributes.
    The functions \texttt{AIC()} and \texttt{BIC()} of \texttt{\{stats\}}
    therefore apply directly, and the package no longer defines its own
    \texttt{AIC} and \texttt{BIC} methods. Section~8.1 applies
    \texttt{logLik()}, which evaluates the log-likelihood at the posterior means
    of the population parameters.
    \item Comment: The harmonic mean estimator has been pointed out to be
    problematic for estimating the marginal likelihood. This aspect should be
    also discussed and covered.
    \item[] Response: We agree and have removed the harmonic mean estimator from
    the package. As \textcite{FruehwirthSchnatter2004} points out, it is a
    reciprocal importance sampling estimator with the prior as importance
    density and therefore unreliable whenever the prior has heavier tails than
    the posterior. The marginal likelihood is now estimated by bridge sampling
    \parencite{Meng1996}, which is much more robust to the tail behavior of the
    importance density \parencite{FruehwirthSchnatter2004}. Section~8.3 applies
    it to the train models, with a result in line with the PSIS-LOO comparison
    of Section~8.2. For models without latent classes, the package uses
    \texttt{\{bridgesampling\}} \parencite{Gronau2020} with a normal importance
    density fitted to the population parameters on an unconstrained scale. Since
    a single normal importance density is unsuitable for the multimodal
    posterior of a mixture \parencite{FruehwirthSchnatter2004}, the importance
    density of a model with a fixed number of up to six latent classes averages
    a normal density, fitted to the relabeled draws, over all permutations of
    the class labels, following the idea of full permutation bridge sampling
    \parencite{Celeux2019,FruehwirthSchnatter2019b}.
\end{enumerate}

\noindent \textbf{Changes from \texttt{\{RprobitB\}} 1.1.0 to
\texttt{\{RprobitB\}} 2.0.1:}

\begin{itemize}
\item The single function \texttt{fit()} replaces the multi-step workflow. It
accepts data in wide or long layout and simulates data if \texttt{data = NULL}.
Data handling and simulation are provided by the companion package
\texttt{\{choicedata\}} \parencite{Oelschlaeger2026a}.
\item Random effects are specified by name through \texttt{random\_effects},
with correlated or uncorrelated normal, log-normal, and negative log-normal
mixing distributions. The argument \texttt{latent\_class\_effects} names the
effects that differ between the latent classes.
\item With \texttt{class\_update}, the number of classes is fixed or adapted
during the warmup by the weight-based heuristic, or a sparse finite or a
Dirichlet process mixture infers the number of occupied classes. The function
\texttt{latent\_class\_diagnostics()} summarizes its posterior and the class
memberships.
\item \texttt{interpret()} computes compensations between effects, such as the
willingness to pay, and marginal effects, each with posterior uncertainty.
\item \texttt{predict()} accounts for posterior uncertainty and provides
conditional predictions given the choice history of a decider.
\texttt{residuals()} is new, and \texttt{update()} is adapted to \texttt{fit()}.
\item The retained draws are stored as a \texttt{draws\_array} object of
\texttt{\{posterior\}}. The function \texttt{summary()} reports the convergence
diagnostics of \textcite{Vehtari2021}, \texttt{plot()} uses
\texttt{\{bayesplot\}}, and chains can run in parallel via \texttt{\{future\}}.
\item The \proglang{C++} implementation of the Gibbs sampler is more efficient.
For instance, the models of Sections~3.5 and 5.1 run about two times faster than
with version 1.1.0.
\end{itemize}

\noindent \textbf{References:}

\printbibliography[heading=none]

\end{document}
